% Source: https://artofproblemsolving.com/community/c3h2192883_solution_without_calculus
% Stable identity: aops-topic-2192883-q01
Put the road on the $x$-axis, Alex at $A=(0,10)$ and the camp at $B=(45,10)$ (both are $10$ km from the road and on the same side, so $AB$ is parallel to the road and has length $45$).

\textbf{Routes that avoid the road.} The shortest such route is the segment $AB$, which takes $\dfrac{45}{30}=1.5$ hours.

\textbf{Routes that use the road.} Let $P=(p,0)$ be the first and $Q=(q,0)$ the last point of the road that the route reaches. The parts before $P$ and after $Q$ are in the desert, so they take at least $\dfrac{AP}{30}$ and $\dfrac{QB}{30}$ hours, and the part from $P$ to $Q$ takes at least $\dfrac{|q-p|}{50}$ hours, because nowhere is the speed more than $50$. Put $u=p$ and $w=45-q$; then $|q-p|\ge45-u-w$, and the Pythagorean theorem gives the time
\[T\ \ge\ \dfrac{\sqrt{100+u^2}}{30}+\dfrac{\sqrt{100+w^2}}{30}+\dfrac{45-u-w}{50}
=\dfrac{5\sqrt{100+u^2}-3u}{150}+\dfrac{5\sqrt{100+w^2}-3w}{150}+\dfrac{9}{10}.\]
For every real $u$,
\[25\left(100+u^2\right)-(40+3u)^2=16u^2-240u+900=(4u-30)^2\ge0,\]
so $5\sqrt{100+u^2}\ge|40+3u|\ge40+3u$, that is, $5\sqrt{100+u^2}-3u\ge40$, with equality exactly when $u=7.5$. The same holds for $w$. Hence
\[T\ \ge\ \dfrac{40}{150}+\dfrac{40}{150}+\dfrac{9}{10}=\dfrac{8}{15}+\dfrac{9}{10}=\dfrac{43}{30}\text{ hours}.\]
This bound is attained: drive from $A$ to $P=(7.5,0)$, along the road to $Q=(37.5,0)$, and on to $B$.

\begin{asy}
usepackage("amsmath");
size(9cm);
defaultpen(fontsize(10pt));
real s=0.16;
pair A=(0,10)*s, B=(45,10)*s, P=(7.5,0)*s, Q=(37.5,0)*s, A0=(0,0), B0=(45,0)*s;
draw((-3,0)*s--(48,0)*s,gray(0.5)+linewidth(2.5));
label("road",(48,0)*s,E,gray(0.4)+fontsize(9pt));
draw(A--B,heavygreen+dashed);
label("$45$ (desert only)",(A+B)/2,N,heavygreen+fontsize(9pt));
draw(A--A0,dotted); draw(B--B0,dotted);
label("$10$",(A+A0)/2,W,fontsize(9pt)); label("$10$",(B+B0)/2,E,fontsize(9pt));
draw(A--P--Q--B,blue+linewidth(1.1));
label("$12.5$",(A+P)/2,E,blue+fontsize(9pt)); label("$12.5$",(Q+B)/2,W,blue+fontsize(9pt));
label("$30$",(P+Q)/2,N,blue+fontsize(9pt));
label("$7.5$",(A0+P)/2,S,fontsize(8pt)); label("$7.5$",(Q+B0)/2,S,fontsize(8pt));
dot(A); dot(B); dot(P); dot(Q);
label("$A$ (Alex)",A,NW); label("$B$ (camp)",B,NE); label("$P$",P,S); label("$Q$",Q,S);
label("\mbox{}",(min(currentpicture,user=true)+max(currentpicture,user=true))/2,black); // keeps the text after the figure black when LaTeX typesets the labels inline
\end{asy}

Each desert leg is $\sqrt{10^2+7.5^2}=12.5$ km and takes $\dfrac{12.5}{30}=\dfrac{5}{12}$ hour, and the $30$ km on the road take $\dfrac{30}{50}=\dfrac{3}{5}$ hour, a total of $\dfrac{5}{6}+\dfrac{3}{5}=\dfrac{43}{30}$ hour. Since $\dfrac{43}{30}<\dfrac{45}{30}$, this beats the straight route.

The fastest trip takes $\boxed{\dfrac{43}{30}\text{ hours}=1\text{ hour }26\text{ minutes}}$.

\emph{Remark.} At the optimum each desert leg makes an angle $\theta$ with the perpendicular to the road where $\sin\theta=\dfrac{7.5}{12.5}=\dfrac{3}{5}=\dfrac{30}{50}$, the ratio of the two speeds (as in Snell's law of refraction).
